\chapter{Background}\label{chapter:background}

\section{Transformer Language Models}\label{sec:bg-transformer}

A transformer language model is a stack of layers~\cite{vaswani2017attention}.
Each layer applies two sub-blocks in sequence: a \emph{self-attention} block, in which each position attends to itself and earlier positions in the sequence, and a \emph{feed-forward network} (FFN), which transforms each token's representation independently.
GPT-OSS-20B has 24 such layers; GPT-OSS-120B has 36~\cite{openai2025gptoss}.

\begin{figure}[htbp]
\centering
\resizebox{0.55\textwidth}{!}{% Figure: high-level decoder-only transformer architecture, focused on one block.
% Goes in Chapter 2, Section 2.1 (Background / Transformer LLM Inference).
%
% Conventions:
%   * Pre-norm decoder block (RMSNorm before attention/FFN), as used by GPT-OSS
%   * Self-attention sub-block: black-boxed; refers to Eq. (attention)
%   * MoE FFN sub-block: highlighted in orange (matches the SSD-storage color
%     used for expert weights elsewhere); refers to Fig. (moe-ffn-dataflow)
%   * Residual connections drawn as green right-side loops
%   * x L badge indicates block stacking (L = 24 for GPT-OSS-20B, 36 for 120B)

\begin{tikzpicture}[
    font=\sffamily\small,
    >={Latex[length=2.5mm]},
    block/.style={
        draw, rectangle, rounded corners=1pt,
        minimum width=30mm, minimum height=8mm,
        align=center, font=\sffamily\small,
        fill=TUMAccentLightBlue!25,
    },
    norm/.style={block, fill=TUMLightGray!30, minimum width=24mm},
    moeblock/.style={
        block, fill=TUMAccentOrange!25, draw=TUMAccentOrange!75!black,
        font=\sffamily\bfseries\small,
    },
    plus/.style={
        draw, circle, fill=white, inner sep=0pt,
        minimum size=4mm, font=\small,
    },
    iolabel/.style={font=\sffamily\small\itshape, text=TUMDarkGray},
    arrow/.style={->, thick},
    residualstyle/.style={->, thick, draw=TUMAccentGreen!60!black},
    lbadge/.style={
        draw=TUMSecondaryBlue, fill=TUMAccentLightBlue!15,
        rounded corners=2pt, font=\sffamily\bfseries\small,
        text=TUMSecondaryBlue, inner sep=4pt,
    },
    xref/.style={font=\scriptsize\sffamily, anchor=west, inner sep=2pt},
]

  % --- Compressed input ---
  \node[iolabel, anchor=south] (inp) at (0, -1.0) {Embedding(tokens)};

  % --- One transformer block, bottom-up ---
  \node[norm]      (norm1) at (0, 0.5) {RMSNorm};
  \node[block]     (attn)  at (0, 1.9) {Self-Attention};
  \node[plus]      (plus1) at (0, 3.3) {$+$};
  \node[norm]      (norm2) at (0, 4.7) {RMSNorm};
  \node[moeblock]  (ffn)   at (0, 6.1) {MoE FFN};
  \node[plus]      (plus2) at (0, 7.5) {$+$};

  % --- Compressed output ---
  \node[iolabel, anchor=north] (out) at (0, 9.2) {RMSNorm $\to$ LM head $\to$ logits};

  % --- Vertical flow ---
  \draw[arrow] (inp)   -- (norm1);
  \draw[arrow] (norm1) -- (attn);
  \draw[arrow] (attn)  -- (plus1);
  \draw[arrow] (plus1) -- (norm2);
  \draw[arrow] (norm2) -- (ffn);
  \draw[arrow] (ffn)   -- (plus2);
  \draw[arrow] (plus2) -- (out);

  % --- Residual loops (right side) ---
  \draw[residualstyle] (0, -0.3) -- (2.4, -0.3) -- (2.4, 3.3) -- (plus1);
  \draw[residualstyle] (0,  3.9) -- (2.4,  3.9) -- (2.4, 7.5) -- (plus2);

  % --- Block boundary ---
  \begin{scope}[on background layer]
    \node[
      draw=TUMGray, dashed, rounded corners=2pt,
      fit=(norm1) (plus2),
      inner xsep=8mm, inner ysep=4mm,
    ] (blockbox) {};
  \end{scope}

  % --- x L badge (left of block boundary) ---
  \node[lbadge, anchor=east] at ($(blockbox.west) + (-0.2, 0)$) {$\times\, L$};

  % --- Cross-references on the right ---
  \node[xref, text=TUMAccentOrange!85!black] at (3.0, 6.1) {See Fig.~\ref{fig:moe-ffn-dataflow}};

  % --- Residual color key ---
  \node[xref, text=TUMAccentGreen!60!black, anchor=west] at (3.0, 4.7) {residual};

\end{tikzpicture}
}
\caption[Decoder-only transformer block]{One transformer layer. Self-attention and the feed-forward network are applied in sequence with residual connections; the model is a stack of $L$ such layers ($L=24$ for GPT-OSS-20B, $L=36$ for GPT-OSS-120B) followed by a final norm and the output projection.}
\label{fig:transformer-block}
\end{figure}

Generating text happens in two phases.
The first, called \emph{prefill}, processes the entire prompt.
The second, called \emph{decode}, produces one new token at a time, each appended to the input for the next step.
Our throughput metric counts tokens generated per second during decode.

To avoid recomputing attention over all prior tokens at every decode step, each layer caches its key and value tensors.
We call this the \emph{KV cache}.
GPT-OSS interleaves layers using a 128-token sliding window with layers using full attention~\cite{openai2025gptoss}, but the KV cache is allocated up front for the full configured 131072-token context.
It therefore reserves 3.02~GiB of RAM on GPT-OSS-20B and 4.53~GiB on GPT-OSS-120B regardless of how many tokens are actually generated.


\section{Mixture-of-Experts}\label{sec:bg-moe}

In a Mixture-of-Experts layer, the FFN is replaced by many small subnetworks, called \emph{experts}~\cite{shazeer2017moe}.
For each generated token, a small \emph{router} selects only a few of them, and the rest are not read.
GPT-OSS-20B has 32 experts per layer and selects 4; GPT-OSS-120B has 128 experts per layer and also selects 4~\cite{openai2025gptoss}.

Each expert is three weight matrices, called the \emph{up}, \emph{gate}, and \emph{down} projections.
We call each of these matrices a \emph{slice}.
A slice is 4.20~MiB on disk on both GPT-OSS models (4{,}406{,}400 bytes).
One MoE layer therefore reads $4 \times 3 = 12$ slices per generated token.
Across the 24 layers of GPT-OSS-20B that is 288 slices per token, 1.18~GiB of weight data.
On GPT-OSS-120B it is $4 \times 3 \times 36 = 432$ slices, 1.77~GiB.

\begin{figure}[htbp]
\centering
\resizebox{0.95\textwidth}{!}{% Figure: MoE-FFN dataflow for one decode token, GPT-OSS-MoE / llama.cpp.
% Goes in Chapter 2 (background).
%
% Architecture-only view: what is computed, with what shapes, in what order,
% for one decode token at one MoE layer. Independent of weight residency.
%
% Conventions (also used by the timing variants in Chapter 7):
%   * activations  : light blue tiles
%   * expert weights (live on SSD, must be loaded): red tiles
%   * compute ops  : green rounded boxes
%   * routing      : violet rounded box
%
% Anchored facts (verified empirically against
% tensor-tracing/desktopui/data/memory-map.json):
%   d := n_embd = ffn_dim = 2880    (GPT-OSS-20B and -120B)
%   k := n_expert_used = 4 (top-k routing)
%   n_e := n_expert = 32 (20B), 128 (120B)
%   per-expert per-projection slice = 2880 * 2880 weights, MXFP4 = 4.20 MiB on disk
%   SwiGLU-OAI: alpha = 1.702, limit = 7.0    (llama-graph.cpp:1764-1765)

\begin{tikzpicture}[
  x=1mm, y=1mm,
  font=\footnotesize,
  every node/.style={inner sep=1pt, outer sep=0pt},
  act/.style    ={draw=blue!55!black,   fill=blue!10,    line width=0.4pt, rounded corners=0.6pt, align=center, font=\scriptsize},
  weight/.style ={draw=red!55!black,    fill=red!12,     line width=0.4pt, rounded corners=0.6pt, align=center, font=\scriptsize},
  op/.style     ={draw=green!55!black,  fill=green!12,   line width=0.5pt, rounded corners=1.2pt, align=center, font=\scriptsize},
  router/.style ={draw=violet!55!black, fill=violet!12,  line width=0.5pt, rounded corners=1.2pt, align=center, font=\scriptsize},
  flow/.style       ={->, >={Stealth[length=1.5mm]}, line width=0.4pt, draw=black!75},
  smallflow/.style  ={->, >={Stealth[length=1.0mm]}, line width=0.3pt, draw=black!60},
  panel/.style      ={draw=black!35, dashed, line width=0.35pt, rounded corners=1.2pt},
  expicon/.style    ={draw=black!50, fill=gray!8, line width=0.3pt, rounded corners=0.6pt, align=center, font=\scriptsize},
]

  % ============================================================
  % FIGURE TITLE: name the operation
  % Centred on image center x=60 (where y[1,d] also lives, at the bottom).
  % ============================================================
  \node[anchor=south, font=\small\bfseries] at (60, 136)
        {Mixture-of-Experts FFN block};
  \node[anchor=south, font=\scriptsize\itshape, text=black!65] at (60, 132)
        {one decode token, one layer; $k\!=\!4$ of $n_e\!=\!32$ experts activated};

  % ============================================================
  % HEADER ROW: input vector x (centred on x=60 so it sits directly above y)
  % ============================================================
  \node[act, minimum width=22mm, minimum height=4mm] (x) at (60, 124)
        {$\mathbf{x}\,[1,d]$};

  % ============================================================
  % ROUTER block (centred on x=60)
  % ============================================================
  \node[router, minimum width=46mm, minimum height=10mm] (router) at (60, 114)
        {\textbf{Router}\\[0.2ex]\scriptsize linear\,$\mathbf{x}\!\cdot\!\mathbf{W}_{r}$
         $\to$\,softmax $\to$ top-$k$};
  \draw[flow] (x) -- (router);

  \node[weight, minimum width=18mm, minimum height=4mm, font=\tiny] (Wr) at (95, 114)
        {$\mathbf{W}_{r}\,[d,n_e]$};
  \draw[smallflow] (Wr) -- (router);

  % Router outputs: placed symmetrically around x=60 (centres at 38 and 82
  % so the section is balanced on the figure's vertical centerline).
  \node[act, minimum width=34mm, minimum height=6mm] (sel) at (38, 102)
        {selected\_experts\,$[k]$\\[-0.4ex]\scriptsize\texttt{[5,\,12,\,23,\,31]}};
  \node[act, minimum width=40mm, minimum height=6mm] (rw) at (82, 102)
        {router\_weights\,$[k]$\\[-0.4ex]\scriptsize\texttt{[0.31,\,0.27,\,0.24,\,0.18]}};

  \draw[flow] (router.south) -| (sel.north);
  \draw[flow] (router.south) -| (rw.north);

  % ============================================================
  % CALLOUT: expert e = 0 in detail.
  % Layout intent (mirrors what each operation actually consumes):
  %   y=82  Wg | x | Wu          three inputs to the gate/up step
  %   y=70  gate matmul | up matmul        parallel matmuls
  %   y=62  gate_0 | up_0                 their two outputs
  %   y=52  swiglu(gate_0, up_0)          combiner
  %   y=42  act_0 | Wd                    two inputs to the down step
  %   y=30  down matmul                   single matmul
  %   y=22  down_0                        per-expert contribution
  % ============================================================
  \draw[panel] (-2, 16) rectangle (60, 88);
  % Title placed ABOVE the panel so it labels the box rather than living inside it.
  \node[anchor=south west, font=\scriptsize\itshape, text=black!65]
        at (-2, 89) {Detail of one expert (representative; \,$e\!=\!0$)};

  % --- top row: parallel inputs to gate/up ---
  \node[weight, minimum width=18mm, minimum height=4mm] (Wg0) at (14, 82)
        {$\mathbf{W}_{g}^{[0]}\,[d,d]$};
  \node[act,    minimum width=10mm, minimum height=4mm] (x0)  at (29, 82)
        {$\mathbf{x}\,[1,d]$};
  \node[weight, minimum width=18mm, minimum height=4mm] (Wu0) at (44, 82)
        {$\mathbf{W}_{u}^{[0]}\,[d,d]$};

  % --- gate / up matmul nodes (green box = "compute op", no need to label "matmul") ---
  \node[op, minimum width=18mm, minimum height=7mm, font=\footnotesize] (gate0) at (14, 70)
        {$\mathbf{x}\!\cdot\!\mathbf{W}_{g}^{[0]\!\top}$};
  \node[op, minimum width=18mm, minimum height=7mm, font=\footnotesize] (up0)   at (44, 70)
        {$\mathbf{x}\!\cdot\!\mathbf{W}_{u}^{[0]\!\top}$};

  % straight feeders Wg -> gate, Wu -> up
  \draw[smallflow] (Wg0) -- (gate0);
  \draw[smallflow] (Wu0) -- (up0);
  % x feeds BOTH gate and up (this is the parallelism the figure makes visible)
  \draw[smallflow] (x0.south) to[out=-90, in=60]  (gate0.north east);
  \draw[smallflow] (x0.south) to[out=-90, in=120] (up0.north west);

  % --- gate / up output activations ---
  \node[act, minimum width=18mm, minimum height=4mm] (g0) at (14, 62) {gate$_{0}\,[1,d]$};
  \node[act, minimum width=18mm, minimum height=4mm] (u0) at (44, 62) {up$_{0}\,[1,d]$};
  \draw[smallflow] (gate0) -- (g0);
  \draw[smallflow] (up0)   -- (u0);

  % --- swiglu combiner (consumes both) ---
  \node[op, minimum width=54mm, minimum height=8mm, font=\footnotesize] (sw0) at (29, 52)
        {$\mathrm{swiglu}(\,\mathrm{gate}_{0}[1,d],\;\mathrm{up}_{0}[1,d]\,)$};
  \draw[smallflow] (g0) -- ([xshift=-7mm]sw0.north);
  \draw[smallflow] (u0) -- ([xshift= 7mm]sw0.north);

  % --- act_0 (left) and Wd (right) on the SAME row, both feeding the down matmul ---
  \node[act,    minimum width=18mm, minimum height=4mm] (a0)  at (15, 42)
        {act$_{0}\,[1,d]$};
  \node[weight, minimum width=18mm, minimum height=4mm] (Wd0) at (43, 42)
        {$\mathbf{W}_{d}^{[0]}\,[d,d]$};
  % swiglu output funnels into act_0 (curving left so it ends pointing into a0)
  \draw[smallflow] (sw0.south) to[out=-90, in=80] (a0.north);

  % --- down matmul (single, takes act_0 and Wd as inputs) ---
  \node[op, minimum width=22mm, minimum height=7mm, font=\footnotesize] (down0) at (29, 30)
        {$\mathrm{act}_{0}\!\cdot\!\mathbf{W}_{d}^{[0]\!\top}$};
  \draw[smallflow] (a0.south)  to[out=-90, in=120] (down0.north west);
  \draw[smallflow] (Wd0.south) to[out=-90, in=60]  (down0.north east);

  % --- down output tile (the per-expert contribution from e=0) ---
  \node[act, minimum width=22mm, minimum height=4mm] (do0) at (29, 22)
        {down$_{0}\,[1,d]$};
  \draw[smallflow] (down0) -- (do0);

  % ============================================================
  % EXPERT ICONS for e = 1, 2, 3 to the right of the callout.
  % Each icon has the SAME outer height as the dashed callout (y in [16, 88])
  % and is wide enough that "down_i [1,d]" fits without clipping.
  % Element ordering inside an icon mirrors the detail panel exactly.
  % ============================================================
  \begin{scope}[shift={(72, 0)}]
    % Single header above the whole 3-icon row (placed once, not per-box)
    \node[anchor=south west, font=\scriptsize\itshape, text=black!65] at (-10, 89)
          {Other experts (same structure, same compute):};

    \foreach \i/\eid in {1/12, 2/23, 3/31} {
      \pgfmathsetmacro{\xi}{(\i-1)*22}
      % Outer icon box: height matches callout (72 mm), width 20 mm
      \node[expicon, minimum width=20mm, minimum height=72mm, anchor=north]
            (e\i) at (\xi, 88) {};
      % Header on a SINGLE row: "e=i (eid=N)"
      \node[anchor=north, font=\scriptsize, text=black!85]
            at (\xi, 86) {$e\!=\!\i$\;\textit{(eid=\eid)}};
      % Vertical stack mirroring the callout: Wg, Wu, gate, up, swiglu,
      % Wd, down, then the down_i output tile separated at the bottom.
      \node[weight, minimum width=15mm, minimum height=4mm, font=\scriptsize] at (\xi, 78) {$\mathbf{W}_{g}^{[\i]}$};
      \node[weight, minimum width=15mm, minimum height=4mm, font=\scriptsize] at (\xi, 72) {$\mathbf{W}_{u}^{[\i]}$};
      \node[op,     minimum width=15mm, minimum height=4mm, font=\scriptsize] at (\xi, 66) {gate};
      \node[op,     minimum width=15mm, minimum height=4mm, font=\scriptsize] at (\xi, 60) {up};
      \node[op,     minimum width=15mm, minimum height=4mm, font=\scriptsize] at (\xi, 54) {swiglu};
      \node[weight, minimum width=15mm, minimum height=4mm, font=\scriptsize] at (\xi, 46) {$\mathbf{W}_{d}^{[\i]}$};
      \node[op,     minimum width=15mm, minimum height=4mm, font=\scriptsize] at (\xi, 40) {down};
      \node[act,    minimum width=17mm, minimum height=4mm, font=\scriptsize, inner sep=0.7pt]
            (di\i_local) at (\xi, 28) {down$_{\i}\,[1,d]$};
    }
  \end{scope}

  % (Routing dashed arrows removed: selected_experts feeding the four expert
  % blocks and router_weights feeding the weighted sum are obvious from the
  % spatial layout; the arrows added clutter without information.)

  % ============================================================
  % WEIGHTED SUM BAND
  % ============================================================
  % Down output tiles in a row above the weighted-sum op
  % (down_0 already drawn inside the callout at (20, 22). The icons drew
  % down_1..down_3 inside their boxes at local x=70..100, y=41.5.)

  \node[op, minimum width=120mm, minimum height=10mm, font=\footnotesize] (sum) at (60, 9)
        {\textbf{Weighted sum:}\;\;
         $\mathbf{y} \;=\; \sum_{e=0}^{k-1}\;\mathrm{router\_weights}[e]\,\cdot\,\mathrm{down}_{e}
         \quad\to\quad [1,d]$};

  % connectors from each down_i into the sum (vertical drops onto sum.north)
  \draw[smallflow] (do0.south) -- ([xshift=-31mm]sum.north);
  \draw[smallflow] (72,  26) -- ([xshift=12mm]sum.north);
  \draw[smallflow] (94,  26) -- ([xshift=34mm]sum.north);
  \draw[smallflow] (116, 26) -- ([xshift=56mm]sum.north);

  % (router_weights -> Weighted sum dashed arrow removed; the formula inside
  % the sum block already states the dependence explicitly.)

  % ============================================================
  % FINAL OUTPUT (lowered slightly so the sum -> y arrow has the same visible
  % length as the x -> Router arrow at the top of the figure)
  % ============================================================
  \node[act, minimum width=22mm, minimum height=4mm] (y) at (60, -1)
        {$\mathbf{y}\,[1,d]$};
  \draw[flow] (sum) -- (y);
  \node[anchor=west, font=\scriptsize\itshape, text=black!65] at (75, -1)
        {to residual\,+\,RMSNorm\,+\,next layer};

  % ============================================================
  % LEGEND (right column, well separated)
  % ============================================================
  \begin{scope}[shift={(132, 110)}, font=\scriptsize]
    \node[anchor=west, font=\scriptsize\bfseries] at (0, 6) {Legend};
    \node[act,    minimum width=8mm, minimum height=3mm] at (4, 0)   {};
      \node[anchor=west] at (8.5, 0)   {activation tile};
    \node[weight, minimum width=8mm, minimum height=3mm] at (4, -5)  {};
      \node[anchor=west] at (8.5, -5)  {expert weight (on SSD)};
    \node[op,     minimum width=8mm, minimum height=3mm] at (4, -10) {};
      \node[anchor=west] at (8.5, -10) {compute op};
    \node[router, minimum width=8mm, minimum height=3mm] at (4, -15) {};
      \node[anchor=west] at (8.5, -15) {routing};

    % Footnote, well below the swatches so no overlap
    \node[anchor=north west, font=\tiny\itshape, text=black!65,
          align=left, text width=42mm] at (0, -22)
         {Red tiles live on SSD; how and when each is loaded is the subject
          of \autoref{chapter:implementation}.};
  \end{scope}

\end{tikzpicture}
}
\caption[Mixture-of-Experts layer for one generated token]{One MoE layer applied to one generated token. The router selects 4 of the 32 experts (highlighted in red); the remaining 28 are not read. Each selected expert applies its three weight matrices (up, gate, down), and the four expert outputs are combined into the layer output.}
\label{fig:moe-ffn-dataflow}
\end{figure}


\section{How llama.cpp Loads Weights}\label{sec:bg-llamacpp}

llama.cpp~\cite{ggerganov2023llamacpp} stores model files in GGUF, a single-file binary format whose header lists each tensor's name, shape, type, and offset, followed by the tensor data~\cite{gguf_spec}.

The default loader maps the file into the process's virtual address space with \texttt{mmap}~\cite{mmap_man}, so tensor bytes are read by dereferencing pointers.
A first access to a not-yet-resident page triggers a \emph{page fault}: the kernel suspends the thread, reads the page (and a small readahead window around it, on the order of tens of KiB) from disk into the \emph{page cache}, maps it at the faulting address, and resumes the thread.
Subsequent accesses to the same page return without entering the kernel.
The page cache lives in kernel-managed memory, and the kernel chooses which pages to evict when other allocations need the space.


\section{Asynchronous Disk I/O on Linux}\label{sec:bg-iouring}

\texttt{io\_uring} is a Linux kernel interface for asynchronous I/O~\cite{axboe2019iouring}.
The application and the kernel share submission and completion queues in memory: a \emph{submission queue} (SQ), into which the application writes read requests, and a \emph{completion queue} (CQ), into which the kernel writes results.
A single \texttt{io\_uring\_enter} system call notifies the kernel that new submission entries are pending. The kernel reads them in place from shared memory and dispatches the I/O without blocking the calling thread.
Each completion entry carries a user-supplied tag that identifies the request it satisfies, so the application can match completions to requests and drain the queue at its own pace.

\texttt{O\_DIRECT} is an \texttt{open(2)} flag that tells the kernel to bypass the page cache for that file~\cite{open_man}.
On a read, the kernel programs the device to DMA the data directly into the application-supplied buffer with no copy through the page cache.
The cost is a strict alignment contract: file offset, read length, and destination buffer address must all be multiples of the device's logical block size, which is 512 bytes on the NVMe used here~\cite{nvme_spec}.
The two features are orthogonal: \texttt{io\_uring} can serve cached or direct I/O, and \texttt{O\_DIRECT} can be used with synchronous \texttt{read(2)} as well.

Combining the two moves both the buffer and the eviction decision into the application.
The submission and completion queues are mmap'd into the application's address space at \texttt{io\_uring} setup time, so the application writes submissions and reads completions without entering the kernel except for the one \texttt{io\_uring\_enter} call per batch.
The data buffer is allocated by the application and is the DMA target (\autoref{fig:iouring-read}).
With no page cache between the device and the buffer, there is no kernel-side cache to evict from; if the application keeps a fixed-size pool of buffer slots, eviction is its responsibility (\autoref{fig:iouring-evict}).

\begin{figure}[htbp]
\centering
\resizebox{0.95\textwidth}{!}{% Figure 2.3: read flow under io_uring + O_DIRECT.
% Goes in Chapter 2, Section 2.4. Companion figure (eviction) lives in
% figures/io_uring_evict.tex with the same column layout.
%
% Three columns, identical placement to the eviction figure:
%   Application (left, dashed): SQ ring, CQ ring, buffer of N=5 slots
%   Kernel       (mid, solid) : single rounded box
%   NVMe SSD    (right, solid): single rounded box
%
% Three arrows distinguished by what they actually carry:
%   (orange, syscall)       app -> kernel: io_uring_enter notifies the kernel
%                           that new SQEs are queued. SQEs themselves stay in
%                           the SQ ring, which the kernel reads in place.
%   (green, metadata)       kernel -> CQ: the kernel writes one CQE per read.
%                           A CQE is a (tag, status) record, not data.
%   (blue, data)            NVMe -> free buffer slot: O_DIRECT DMA. This is
%                           the only arrow on this figure that moves bytes
%                           of file content.
%
% Dashed gray "buffer ptr" link from SQE to the destination buffer slot
% answers "how does io_uring know where to write": each SQE carries a
% pointer to the target buffer slot.
%
% A small italic note under the CQ row reinforces that CQEs carry only
% tag + status (no data flows through the CQ).

\begin{tikzpicture}[
    font=\sffamily\small,
    >={Latex[length=2.5mm]},
    sqe/.style={draw=TUMBlue, fill=TUMBlue!15, minimum width=8mm, minimum height=6mm, font=\footnotesize, inner sep=1pt},
    cqe/.style={draw=TUMAccentGreen!80!black, fill=TUMAccentGreen!25, minimum width=8mm, minimum height=6mm, font=\footnotesize, inner sep=1pt},
    bufslot/.style={draw=TUMSecondaryBlue, fill=TUMAccentLightBlue!40, minimum width=10mm, minimum height=6mm, font=\footnotesize, inner sep=1pt},
    bigbox/.style={draw=TUMGray, rounded corners=2pt, minimum width=2.0cm, font=\small\bfseries},
    arrow/.style={->, thick},
    arrowlbl/.style={font=\scriptsize, sloped=false},
    annot/.style={font=\scriptsize\itshape, text=TUMDarkGray},
]

  % --- Application column ---
  \node[font=\small\bfseries, anchor=west] (applabel) at (0.2, 5.0) {Application};

  \node[anchor=east, font=\footnotesize] (sqlabel) at (0.7, 4.1) {SQ};
  \node[sqe, anchor=west] (s0) at (0.8, 4.1) {S\textsubscript{0}};
  \node[sqe, anchor=west] (s1) at (s0.east) {S\textsubscript{1}};
  \node[sqe, anchor=west] (s2) at (s1.east) {S\textsubscript{2}};
  \node[sqe, anchor=west] (s3) at (s2.east) {S\textsubscript{3}};
  \node[sqe, anchor=west, draw=TUMBlue!50, fill=white] (s4) at (s3.east) {};

  \node[anchor=east, font=\footnotesize] (cqlabel) at (0.7, 3.0) {CQ};
  \node[cqe, anchor=west] (c0) at (0.8, 3.0) {C\textsubscript{0}};
  \node[cqe, anchor=west] (c1) at (c0.east) {C\textsubscript{1}};
  \node[cqe, anchor=west] (c2) at (c1.east) {C\textsubscript{2}};
  \node[cqe, anchor=west] (c3) at (c2.east) {C\textsubscript{3}};
  \node[cqe, anchor=west, draw=TUMAccentGreen!50, fill=white] (c4) at (c3.east) {};

  % CQ contents annotation
  \node[annot, anchor=north west] (cqannot)
      at ($(c0.south west)+(0,-0.5mm)$) {each CQE: (tag, bytes read), no data};

  \node[anchor=east, font=\footnotesize] (buflabel) at (0.7, 1.5) {buffer};
  \node[bufslot, anchor=west] (b0) at (0.8, 1.5) {slice\textsubscript{0}};
  \node[bufslot, anchor=west] (b1) at (b0.east) {slice\textsubscript{1}};
  \node[bufslot, anchor=west] (b2) at (b1.east) {slice\textsubscript{2}};
  \node[bufslot, anchor=west] (b3) at (b2.east) {slice\textsubscript{3}};
  \node[bufslot, anchor=west, draw=TUMSecondaryBlue!60, fill=TUMAccentLightBlue!10] (bfree) at (b3.east) {free};

  \begin{scope}[on background layer]
    \node[draw=TUMGray, dashed, rounded corners=2pt, fit=(applabel) (sqlabel) (s4) (cqlabel) (c4) (cqannot) (buflabel) (bfree), inner sep=4mm] (appbox) {};
  \end{scope}

  % --- Kernel column ---
  \node[bigbox, fill=TUMLightGray!25, minimum height=3.4cm, anchor=center] (kernel) at (9.6, 3.0) {Kernel};

  % --- NVMe column ---
  \node[bigbox, fill=TUMLightGray!50, minimum height=1.4cm, draw=TUMBlack, anchor=center] (nvme) at (12.6, 1.5) {NVMe};

  % --- Active arrows ---
  % (orange) io_uring_enter is a syscall notifying the kernel.
  % SQEs stay in shared memory; the kernel reads them in place.
  \draw[arrow, TUMAccentOrange]
      (s4.east) -- node[arrowlbl, above] {\texttt{io\_uring\_enter} (syscall)}
      (kernel.west |- s2);

  % (green) kernel writes a CQE into shared memory: tag + status only.
  \draw[arrow, TUMAccentGreen!60!black]
      (kernel.west |- c2) -- node[arrowlbl, above] {write CQE}
      (c4.east);

  % (blue, thicker) DMA: NVMe writes file bytes directly into the free
  % buffer slot, bypassing the kernel column. This is the only arrow on
  % this figure that carries data.
  \draw[arrow, TUMSecondaryBlue, line width=1pt]
      (nvme.west) -- node[arrowlbl, above] {data (DMA, \texttt{O\_DIRECT})}
      (bfree.east);

\end{tikzpicture}
}
\caption[Read flow under \texttt{io\_uring} with \texttt{O\_DIRECT}]{Read flow under \texttt{io\_uring} combined with \texttt{O\_DIRECT}. The submission queue, completion queue, and the data buffer all live in the application's address space. Each SQE carries the destination buffer pointer for its read. One \texttt{io\_uring\_enter} system call (orange) notifies the kernel; the kernel reads the SQEs from shared memory, pins the destination pages, and programs the device. The device then DMAs the file bytes into the destination, overwriting whatever was there (blue): this is the only path on the figure that carries data. The kernel writes one CQE per read into the CQ ring (green); a CQE is a (tag, bytes-read) record only, no data. The application reads the CQ in place without another syscall.}
\label{fig:iouring-read}
\end{figure}

\begin{figure}[htbp]
\centering
\resizebox{0.95\textwidth}{!}{% Figure 2.4: eviction step under io_uring + O_DIRECT.
% Same column placement and same building blocks as figure 2.3
% (figures/io_uring_rings.tex) so the two figures contrast cleanly.
%
% On a cache miss with the buffer already full, the application picks a
% victim slot under its replacement policy (LRU, LFU, or LFU-aging in our
% backend). The miss is then served by the same SQ -> io_uring_enter ->
% DMA -> CQ flow shown in figure 2.3, with the freshly freed slot as the
% DMA target. Eviction itself is purely application-internal: no syscall,
% no device traffic.
%
% Active arrow on this figure:
%   (red)    application policy marks slice_3 as the victim slot
% Inactive arrows:
%   (gray, dashed) the io_uring rings and DMA path drawn in light gray to
%                  show the same building blocks are still in place.

\begin{tikzpicture}[
    font=\sffamily\small,
    >={Latex[length=2.5mm]},
    sqe/.style={draw=TUMBlue!40, fill=TUMBlue!8, minimum width=8mm, minimum height=6mm, font=\footnotesize, inner sep=1pt, text=TUMDarkGray},
    cqe/.style={draw=TUMAccentGreen!40, fill=TUMAccentGreen!10, minimum width=8mm, minimum height=6mm, font=\footnotesize, inner sep=1pt, text=TUMDarkGray},
    bufslot/.style={draw=TUMSecondaryBlue, fill=TUMAccentLightBlue!40, minimum width=10mm, minimum height=6mm, font=\footnotesize, inner sep=1pt},
    bigbox/.style={draw=TUMGray, rounded corners=2pt, minimum width=2.0cm, font=\small\bfseries},
    arrow/.style={->, thick},
    arrowlbl/.style={font=\scriptsize, sloped=false},
    inactive/.style={draw=TUMGray!50, dashed, ->, thick},
]

  % --- Application column (dashed border) ---
  \node[font=\small\bfseries, anchor=west] (applabel) at (0.2, 5.0) {Application};

  \node[anchor=east, font=\footnotesize] (sqlabel) at (0.7, 4.1) {SQ};
  \node[sqe, anchor=west] (s0) at (0.8, 4.1) {};
  \node[sqe, anchor=west] (s1) at (s0.east) {};
  \node[sqe, anchor=west] (s2) at (s1.east) {};
  \node[sqe, anchor=west] (s3) at (s2.east) {};
  \node[sqe, anchor=west] (s4) at (s3.east) {};

  \node[anchor=east, font=\footnotesize] (cqlabel) at (0.7, 3.0) {CQ};
  \node[cqe, anchor=west] (c0) at (0.8, 3.0) {};
  \node[cqe, anchor=west] (c1) at (c0.east) {};
  \node[cqe, anchor=west] (c2) at (c1.east) {};
  \node[cqe, anchor=west] (c3) at (c2.east) {};
  \node[cqe, anchor=west] (c4) at (c3.east) {};

  \node[anchor=east, font=\footnotesize] (buflabel) at (0.7, 1.5) {buffer};
  \node[bufslot, anchor=west] (b0) at (0.8, 1.5) {slice\textsubscript{0}};
  \node[bufslot, anchor=west] (b1) at (b0.east) {slice\textsubscript{1}};
  \node[bufslot, anchor=west] (b2) at (b1.east) {slice\textsubscript{2}};
  \node[bufslot, anchor=west,
        draw=red!70!black, fill=red!20, line width=1pt] (bvictim) at (b2.east) {slice\textsubscript{3}};
  \node[bufslot, anchor=west] (b4) at (bvictim.east) {slice\textsubscript{4}};

  \begin{scope}[on background layer]
    \node[draw=TUMGray, dashed, rounded corners=2pt, fit=(applabel) (sqlabel) (s4) (cqlabel) (c4) (buflabel) (b4), inner sep=4mm] (appbox) {};
  \end{scope}

  % --- Kernel column (drawn in light gray to indicate it is not active) ---
  \node[bigbox, fill=TUMLightGray!15, minimum height=3.4cm, draw=TUMGray!50, text=TUMDarkGray, anchor=center] (kernel) at (9.6, 3.0) {Kernel};

  % --- NVMe column (also light gray, not active) ---
  \node[bigbox, fill=TUMLightGray!25, minimum height=1.4cm, draw=TUMGray!50, text=TUMDarkGray, anchor=center] (nvme) at (12.6, 1.5) {NVMe};

  % --- Inactive arrows (dashed, gray) just to show the same diagram skeleton ---
  \draw[inactive] (s4.east) -- (kernel.west |- s2);
  \draw[inactive] (kernel.west |- c2) -- (c4.east);
  \draw[inactive] (nvme.west) -- (b4.east);

  % --- The eviction action (red, intra-application) ---
  % Policy node above the buffer
  \node[draw=red!70!black, rounded corners=2pt, fill=red!8,
        font=\footnotesize, align=center, inner sep=2pt]
        (policy) at (3.6, 0.3) {policy picks victim};
  \draw[arrow, red!70!black] (policy.north) -- (bvictim.south);

  % Annotate that no kernel/device traffic occurs
  \node[font=\scriptsize\itshape, text=TUMDarkGray, align=center]
        at (10.6, 0.3) {(no syscall, no\\device traffic)};

\end{tikzpicture}
}
\caption[Eviction under \texttt{io\_uring} with \texttt{O\_DIRECT}]{Eviction step on a cache miss when the buffer is already full, drawn with the same building blocks as \autoref{fig:iouring-read}. The application's replacement policy picks a victim slot (here \emph{slice\textsubscript{3}}, in red); no system call and no device traffic occur. The freed slot then serves as the DMA target for the read flow of \autoref{fig:iouring-read}.}
\label{fig:iouring-evict}
\end{figure}
